\subsection{Accumulated memory and transport resolvents}
\label{subsec:memoria-resolvente}

The return of the calendar does not eliminate the increments produced during
the traversal. To express this distinction in operator terms, we first construct
an accumulated record of path events and then its
generating series. Memory variables will be reduced on
that transport, while also preserving the source and output reader.

\subsubsection{Oriented events and record update}

Let \((d_t)\) be the sequence of integer codes for direction preserved
by a TPK path. The code \(d_t\) and the cursor's cardinal direction
are distinct coordinates; no identification between their alphabets
is assumed. Number the windows from zero:
\[
 I_m=\{9m+1,\ldots,9m+9\},\qquad m\geq0.
\]
In the first cycle, \(I_m\) is window \(E_{m+1}\) of the
dodecaphase register. Its orientation is
\[
 \sigma=(1,1,1,-1,-1,-1,1,1,1,-1,-1,-1).
\]
For an extended history, \(\sigma_m\in\{-1,1\}\) is the sign
preserved in the corresponding window. The signed event is
\begin{equation}
 a_m=\sigma_m\sum_{t\in I_m}
       \mathbf1_{\{2,4,6,8\}}(d_t),\qquad |a_m|\leq9.
 \label{mem:evento}
\end{equation}
Each summand is incorporated when its transition occurs. This definition
counts events in the trace; it does not select directions by means of a value
that is to be obtained later.

Let \(\mathbf e_0,\ldots,\mathbf e_{11}\) be the basis of \(\ZZ^{12}\)
and let \(S\mathbf e_j=\mathbf e_{j+1\bmod12}\). Write
\(R=S^{-1}\). The accumulated record in fixed coordinates satisfies
\[
z_{m+1}=z_m+a_m\mathbf e_{m\bmod12}.
\]
A prefix without accumulated history begins at \(z_0=0\); at another
origin, \(z_0\) preserves the preceding record.
The frame accompanying the calendar is \(y_m=S^{-m}z_m\).

\begin{proposition}[Update and return with memory]
\label{mem:prop-actualizacion}
The record in the moving frame satisfies
\begin{equation}
 y_{m+1}=Ry_m+\eta_m,\qquad
 \eta_m=a_mR\mathbf e_0.
 \label{mem:actualizacion}
\end{equation}
For every \(n\geq1\),
\begin{equation}
 y_n=R^ny_0+\sum_{j=0}^{n-1}R^{n-1-j}\eta_j.
 \label{mem:iteracion}
\end{equation}
In particular,
\begin{equation}
 y_{12}=y_0+\sum_{j=0}^{11}a_j\mathbf e_j.
 \label{mem:retorno12}
\end{equation}
\end{proposition}
\begin{proof}
Multiplication of the \(z_m\)-update by \(S^{-(m+1)}\) yields
\(y_{m+1}=S^{-1}y_m+a_mS^{-1}\mathbf e_0\), because
\(S^{-m}\mathbf e_{m\bmod12}=\mathbf e_0\).

Formula \eqref{mem:iteracion} follows by induction.
For \(n=12\), \(R^{12}=I\) and
\(R^{11-j}\eta_j=a_jR^{12-j}\mathbf e_0=a_j\mathbf e_j\),
which proves \eqref{mem:retorno12}.
\end{proof}

Thus, twelve windows restore the frame and preserve the event vector.
The count alone does not recover the order of every step
within each window; that order remains in the recorded path.

\Needspace{11\baselineskip}
To specify the balance, define
\[
 D_3=I-S^3,\quad D_4=I-S^4,\quad
 B=D_3^{\mathsf T}D_3+D_4^{\mathsf T}D_4,
\]
\[
 \mathcal M(y)=\tfrac12y^{\mathsf T}By,
 \qquad q(y)=\sum_{j=0}^{11}y_j.
\]
The first is a quadratic form of the record; the second records
the sum of its coordinates. They are not identified here with physical energy or charge.

\begin{proposition}[Balance produced by the event]
\label{mem:prop-balance}
Update \eqref{mem:actualizacion} satisfies
\begin{equation}
 \begin{aligned}
 \mathcal M(y_{m+1})-\mathcal M(y_m)
   &=a_m(By_m)_0+2a_m^2,\\
 q(y_{m+1})-q(y_m)&=a_m.
 \end{aligned}
 \label{mem:balance}
\end{equation}
\end{proposition}
\begin{proof}
Expanding the products gives
\(B=4I-S^3-S^{-3}-S^4-S^{-4}\); hence
\(R^{\mathsf T}BR=B\) and \(B_{00}=4\).
Since \(y_{m+1}=R(y_m+a_m\mathbf e_0)\), the quadratic difference
is \(a_m\mathbf e_0^{\mathsf T}By_m+\tfrac12a_m^2B_{00}\).
The sum of the coordinates is invariant under the permutation \(R\),
which yields the second equality.
\end{proof}

\Needspace{20\baselineskip}
\subsubsection{Generating series and uniform control of memory}

The series preserves every increment, not merely their sum at the end
of a cycle. With a formal variable \(u\), let
\[
 Y(u)=\sum_{m\geq0}y_mu^m,\qquad
 H_\eta(u)=\sum_{m\geq0}\eta_mu^m.
\]

\begin{proposition}[Resolvent and tail bound]
\label{mem:prop-serie}
The formal identity
\begin{equation}
 Y(u)=(I-uR)^{-1}\bigl(y_0+uH_\eta(u)\bigr).
 \label{mem:serie}
\end{equation}
holds. Both series converge for \(|u|<1\). If \(\ell:\RR^{12}\to\RR\)
is linear, \(L=\|\ell\|_{1\to\RR}\), \(M_0=\|y_0\|_1\), and
\(0<r<1\), then, for every integer \(N\geq0\) and \(|u|\leq r\),
\begin{equation}
 \left|\sum_{m>N}\ell(y_m)u^m\right|
 \leq Lr^{N+1}\left(
 \frac{M_0+9(N+1)}{1-r}+\frac{9r}{(1-r)^2}\right).
 \label{mem:cota}
\end{equation}
The derivative series converges uniformly on every disk
\(|u|\leq r<1\).
\end{proposition}
\begin{proof}
Multiplying \eqref{mem:actualizacion} by \(u^{m+1}\) and summing gives
\(Y-y_0=uRY+uH_\eta\). The constant term of \(I-uR\) is
invertible, so its formal inverse is \(\sum_{k\geq0}u^kR^k\).
Moreover, \(R\) preserves the \(\ell^1\)-norm and
\(\|\eta_m\|_1\leq9\), whence
\(\|y_m\|_1\leq M_0+9m\).
For the tail, sum the majorants
\(L(M_0+9m)r^m\), using
\[
 \sum_{m>N}r^m=\frac{r^{N+1}}{1-r},\qquad
 \sum_{m>N}mr^m=r^{N+1}
 \left(\frac{N+1}{1-r}+\frac r{(1-r)^2}\right).
\]
Finally, \(\sum_{m\geq1}Lm(M_0+9m)r^{m-1}<\infty\)
majorizes the derivative series and proves its uniform convergence.
\end{proof}

The variable \(u\) counts windows of nine transitions. Identifying it
with the variable of another reader requires preserving that clock and the map
relating the two readings.

\subsubsection{Elimination of memory: operator, source, and reader}

On a domain where the update is fixed as \(U:X\to X\),
the free vector space \(V=\QQ^{(X)}\) on the states admits the operator
\(T\delta_x=\delta_{U(x)}\). The clock is included in \(x\) if the
transition depends on it. This linear extension preserves distinct states
even when they share an observable phase.

Consider a decomposition \(V=V_0\oplus V_1\), where \(V_1\)
is the auxiliary sector to be eliminated, and write
\[
 T=\begin{pmatrix}A&B_1\\C&D\end{pmatrix},\qquad
 v=\binom{v_0}{v_1},\qquad \ell=(\ell_0,\ell_1).
\]
Here \(A:V_0\to V_0\), \(B_1:V_1\to V_0\),
\(C:V_0\to V_1\), and \(D:V_1\to V_1\).
The following identities are formal; they also hold wherever the
resolvents converge.

\begin{proposition}[Schur reduction with source and reader]
\label{mem:prop-schur}
Let
\[
 D_u=(I-uD)^{-1},\qquad
 \mathscr F(u)=I-uA-u^2B_1D_uC.
\]
The retained component of \((I-uT)^{-1}v\) is
\begin{equation}
 w_0=\mathscr F(u)^{-1}\bigl(v_0+uB_1D_uv_1\bigr),
 \qquad w_1=D_u(v_1+uCw_0).
 \label{mem:schur-fuente}
\end{equation}
Its complete reading satisfies
\begin{equation}
 \begin{split}
 \ell(I-uT)^{-1}v
 ={}&(\ell_0+u\ell_1D_uC)
       \mathscr F(u)^{-1}(v_0+uB_1D_uv_1)\\
    &+\ell_1D_uv_1.
 \end{split}
 \label{mem:schur-lector}
\end{equation}
\end{proposition}
\begin{proof}
The two equations of \((I-uT)w=v\) are
\(w_0=v_0+uAw_0+uB_1w_1\) and
\(w_1=v_1+uCw_0+uDw_1\).
Solving the second and substituting it into the first gives
\eqref{mem:schur-fuente}; applying \(\ell_0\) and \(\ell_1\)
to both components gives \eqref{mem:schur-lector}.
All formal inverses exist because their constant terms
are identities.
\end{proof}

The term \(u^2B_1D_uC=\sum_{k\geq0}u^{k+2}B_1D^kC\)
describes an exit into memory, \(k\) steps within it, and
a return. By contrast, \(u\ell_1D_uC\) reads the memory before that
return. Consequently, reducing only the operator does not necessarily preserve
the observation: the source and reader must also be transformed.

The distinction is visible in the cycle already constructed. If \(P\) retains
\(\mathbf e_0\), then \(PRP=0\), but
\begin{equation}
 \left\langle\mathbf e_0,(I-uR)^{-1}\mathbf e_0\right\rangle
 =\frac1{1-u^{12}},\qquad
 \left\langle\mathbf e_{11},(I-uR)^{-1}\mathbf e_0\right\rangle
 =\frac u{1-u^{12}}.
 \label{mem:ejemplo-ciclo}
\end{equation}
Indeed, \(R^n\mathbf e_0=\mathbf e_{-n\bmod12}\): the
two series collect, respectively, \(n\equiv0\) and
\(n\equiv1\pmod{12}\). The resolvent of the compression is \(I\),
whereas compression of the resolvent preserves the omitted returns.
This example pertains to the twelve-window record, not to the entirety
of the TPK state.

\subsubsection{Composition of segments and coherence of reduction}

Let three composable affine transports be
\(F_i(y)=R_i y+b_i\), with compatible domains and codomains.
The memory produced by the complete traversal is
\begin{equation}
 F_3F_2F_1(y)
 =R_3R_2R_1y+b_3+R_3b_2+R_3R_2b_1.
 \label{mem:cociclo-tres}
\end{equation}
The proof is the successive substitution of the three affine formulas.
Grouping either the first two segments or the last two first produces
the same expression; the factors transport each increment from
the point at which it originated. In \eqref{mem:actualizacion},
\(R_i=R\), and \(b_i\) is the event of its window.

Elimination of several layers preserves analogous coherence.
To make it explicit, now decompose
\(V=V_0\oplus V_1\oplus V_2\) and write
\[
 T=\begin{pmatrix}
 A&B_1&B_2\\C_1&D_{11}&D_{12}\\C_2&D_{21}&D_{22}
 \end{pmatrix},\qquad Q_2=(I-uD_{22})^{-1}.
\]
After eliminating the third component, the four blocks are
\begin{equation}
 \begin{aligned}
 A'&=A+uB_2Q_2C_2,&
 B_1'&=B_1+uB_2Q_2D_{21},\\
 C_1'&=C_1+uD_{12}Q_2C_2,&
 D_{11}'&=D_{11}+uD_{12}Q_2D_{21}.
 \end{aligned}
 \label{mem:schur-intermedio}
\end{equation}

\begin{proposition}[Transitivity of elimination]
If \(D_*=(D_{ij})_{i,j=1}^2\), the complement obtained in two stages
coincides with that obtained by eliminating both memories jointly:
\begin{equation}
 \begin{split}
 I-uA'-u^2B_1'(I-uD_{11}')^{-1}C_1'
 ={}&I-uA\\
 &-u^2(B_1\ B_2)(I-uD_*)^{-1}\binom{C_1}{C_2}.
 \end{split}
 \label{mem:schur-transitivo}
\end{equation}
\end{proposition}
\begin{proof}
Solving the third equation of \((I-uT)w=v\) and substituting it into the
first two produces \eqref{mem:schur-intermedio}. Subsequently eliminating the
second component gives the left-hand side of
\eqref{mem:schur-transitivo}. Solving the two memory equations simultaneously
gives the right-hand side, by Proposition~\ref{mem:prop-schur}.
The formal solution is unique because \(I-uT\) has constant term
\(I\). In both procedures, the source and reader are also transformed
according to \eqref{mem:schur-fuente}--\eqref{mem:schur-lector}.
\end{proof}

Balance \eqref{mem:balance} describes the increments of the record.
The preceding identities allow them to be transported when memory is reused
in the dodecaphase closure; by themselves, they do not determine additional
analytic coefficients.
